\section{Effective unimodality under dilation}\label{sec:dilation}

For a fixed nonempty composition $w=(w_1,\ldots,w_l)$, write
$tw=(tw_1,\ldots,tw_l)$ and $W=\sum_iw_i$.

\begin{theorem}[Dilation and its eventual modes]\label{thm:dilation}
There is a computable integer $T(w)$ such that $F_{tw}$ is unimodal
for every integer $t\ge T(w)$.
If $2\max_iw_i<W$, set $C_t=(tW+l-1)/2$. For all $t\ge T(w)$,
the lower integer $\lfloor C_t\rfloor$ is a mode, and it is either
the unique mode or belongs to exactly two adjacent modes
$\lfloor C_t\rfloor,\lceil C_t\rceil$.
In particular, an integral $C_t$ is the unique mode.
The eventual choice in each parity of $t$ is computable exactly.
There is no restriction on the length or number of descents of $w$.
\end{theorem}

The proof confines all undecided differences to a fixed finite window,
then computes their exact polynomials in $t$. A limiting unimodal
density, without control of these differences, would not suffice.

\begin{lemma}[A universal right-tail bound]\label{lem:righttail}
For a nonempty host of weight $N$ and length $l$, its rank polynomial
is a nonnegative sum of products of interval polynomials, each
symmetric about a point at most $(N+l-1)/2$.
Consequently
\[
 [z^n]F_w\le[z^{n-1}]F_w
       \qquad\left(n\ge\left\lceil\frac{N+l}{2}\right\rceil\right).
\]
\end{lemma}
\begin{proof}
For $l=1$, use $F_w=\Iv0N$. Assume $l\ge2$. Set
$m_i=\max(w_{i+1},\ldots,w_l)$, with $m_l=0$, and
$N_i=w_1+\cdots+w_i$. Group the terms of the rightmost-last-letter
identity according to their host position:
\[
 F_w=1+\sum_{i:w_i>m_i}\Iv{m_i+1}{w_i}F_{w_1\cdots w_{i-1}}.
\]
For $i\ge2$, the induction hypothesis bounds twice the center of
each product in this summand by
\[
 N_{i-1}+i-2+(w_i+m_i+1)=N_i+i-1+m_i\le N+l-1,
\]
since $m_i\le N-N_i$. For $i=1$, the twice-center is
$w_1+m_1+1\le N+l-1$, using $l\ge2$.
The constant term also meets the bound. Products of intervals are
symmetric and unimodal, so each is nonincreasing at every difference
index $n\ge\lceil(N+l)/2\rceil$. Summing proves the assertion.
\end{proof}


Together with Lemma~\ref{lem:half}, the right-tail bound leaves only
the differences strictly after $\lceil N/2\rceil$ and no later than
$\lfloor(N+l-1)/2\rfloor$ undecided. The width of this window depends
on the number of parts, not their sizes.

\begin{lemma}[Multilinear expansion]\label{lem:multilinear}
There are integers $a_I$, independent of $t$, such that
\[
 F_{tw}=\sum_{I\subseteq[l]}a_I\prod_{i\in I}S_{tw_i},
 \qquad a_{[l]}=1,\qquad
 \sum_{|I|=j}a_I=1\quad(0\le j\le l).
\]
One may choose them with $\sum_I|a_I|\le3^l$.
\end{lemma}
\begin{proof}
Fix recursive splits at maxima of $w$, which remain valid under
dilation. Apply the maximum identity, retaining a formal variable for
each position. The resulting polynomial is multilinear and its full
monomial has coefficient one. Setting every variable equal to $s$
gives $1+s+\cdots+s^l$: for arms of lengths $a,b$, the expressions
$1+(s-1)P_X=s^{a+1}$ and $1+(s-1)P_Y=s^{b+1}$ multiply.
This proves the sums by cardinality.
An endpoint maximum gives $P=1+sP_Y$. At an internal maximum,
$\|P\|_1\le\|P_X\|_1+\|P_Y\|_1+2\|P_X\|_1\|P_Y\|_1$.
For $a,b\ge1$, $3^a+3^b\le3^{a+b}$, so induction gives the norm bound.
\end{proof}

For positive widths $v_1,\ldots,v_m$, let
$B_v=\mathbf1_{[0,v_1]}*\cdots*\mathbf1_{[0,v_m]}$.
For $m\ge2$ its continuous representative is
\begin{equation}\label{eq:boxspline}
 B_v(x)=\frac1{(m-1)!}
 \sum_{J\subseteq[m]}(-1)^{|J|}
       \left(x-\sum_{j\in J}v_j\right)_+^{m-1}.
\end{equation}
Repeated integration of interval indicators proves this formula.

\begin{lemma}[Negative central curvature]\label{lem:boxcurvature}
If $m\ge3$ and $2\max_i v_i<V=\sum_i v_i$, then
$B_v''(V/2)$ exists and is strictly negative.
\end{lemma}
\begin{proof}
We use the elementary strict shape of interval convolutions. In
centered coordinates, a convolution with total width $V$ and largest
width $L$ has a plateau of half-width $(2L-V)_+/2$ and is strictly
decreasing from that plateau to its support endpoint. For two
intervals this is the trapezoid formula. Inductively put the largest
interval last, and write $g$ for the centered convolution of the other
intervals, of total width $D=V-L$. Then
\[
 B_v'(V/2+x)=g(x+L/2)-g(x-L/2).
\]
For $(L-D)_+/2<x<(L+D)/2$, the inner point lies inside the support,
and the outer point lies beyond the plateau of $g$. That plateau has
half-width strictly less than $L/2$ when at least two widths remain.
Symmetry and the inductive strict decrease make the displayed
difference negative. At $x=0$ the difference is zero by symmetry.
If $L\ge D$, both terms vanish for $0\le x\le(L-D)/2$.
This proves the asserted strict shape, including strict negativity
of its derivative off the plateau for three or more intervals.

Under $L<D$, the point $L/2$ lies strictly between the plateau and
the support endpoint of $g$. Hence
\[
 B_v''(V/2)=2g'(L/2)<0.
\]
When $g$ is a two-interval convolution the point is in a sloping
open segment of its trapezoid; otherwise $g'$ is continuous there.
\end{proof}

Write $c_i=a_{[l]\setminus\{i\}}$, so $\sum_i c_i=1$.
The top two multilinear degrees have an exact symmetry:
\begin{equation}\label{eq:topsymmetry}
 H_t=\prod_i S_{tw_i}+\sum_i c_i\prod_{j\ne i}S_{tw_j},
 \qquad z^{tW+l-1}H_t(1/z)=H_t(z).
\end{equation}
Indeed, reflection multiplies the full product by $z^{-1}$ and the
$i$th omitted-factor product by $z^{tw_i}$. The differences cancel,
using $z^{tw_i}-1=-(1-z)S_{tw_i}/z$ and $\sum_i c_i=1$.

\begin{lemma}[Exact central polynomials]\label{lem:centralpolynomials}
Assume $2\max_iw_i<W$. Choose a parity $\epsilon\in\{0,1\}$ and
a half-integer $s$ with $0\le s\le l/2+1$ and
$\epsilon W/2+s\in\mathbb Z$. There is an explicitly computable
rational polynomial $P_s(t)$ such that
\[
 \Del F_{tw}(tW/2+s)=P_s(t)
 \quad\text{for }t\ge4l+4,\quad t\equiv\epsilon\pmod2.
\]
For $l\ge4$, its degree is at most $l-3$, and its coefficient of
$t^{l-3}$ is
\begin{equation}\label{eq:centralleading}
 B_w''(W/2)(s-l/2).
\end{equation}
For $l=3$ the polynomial is exactly $-2s+3$.
At $s=l/2$, when this parity exists, write $Q_w=P_{l/2}$.
This polynomial is either zero or has negative leading coefficient;
for $l\ge4$ its degree is at most $l-4$.
\end{lemma}
\begin{proof}
For a product of $m\ge2$ intervals the exact first-difference formula is
\begin{equation}\label{eq:centralbinomial}
 \Del\left(\prod_{i\in I}S_{tw_i}\right)(n)
 =\sum_{J\subseteq I}(-1)^{|J|}
    \binom{n-t\sum_{j\in J}w_j-2}{m-2}.
\end{equation}
A binomial is zero unless its upper argument is at least its
nonnegative lower argument. This follows by expanding the numerator
of $z^m\prod_{i\in I}(1-z^{tw_i})/(1-z)^{m-1}$.
At $n=tW/2+s$, a summand has upper argument
\[
 t\beta+s-2,\qquad \beta=W/2-\sum_{j\in J}w_j.
\]
For $t\ge4l+4$, a positive $\beta$ is at least $1/2$ and makes the
binomial an ordinary polynomial in $t$; a negative $\beta$ is at
most $-1/2$ and makes it zero. If $\beta=0$, retain its constant
binomial if it is nonzero. In that case $W$ is even and the allowed
parity makes $s$ integral. Products with fewer than two factors have
zero difference in this balanced central window. This constructs
$P_s$ exactly from Lemma~\ref{lem:multilinear}.

For $l\ge4$, the leading two terms of the full product's difference
are
\[
 t^{l-2}B_w'(W/2)
 +t^{l-3}\left(s-\frac{l+1}{2}\right)B_w''(W/2).
\]
To verify the second coefficient directly, the shifts in the
binomial product are $2,3,\ldots,l-1$, whose sum is
$(l-2)(l+1)/2$. The first term vanishes by symmetry.
An omitted-factor product contributes leading coefficient
$B_{w\setminus i}'(W/2)$ at degree $l-3$. Differentiating the
convolution with the $i$th interval and using symmetry gives
\[
 B_w''(W/2)=B_{w\setminus i}'(W/2)
              -B_{w\setminus i}'(W/2-w_i)
           =2B_{w\setminus i}'(W/2).
\]
All other products are of degree at most $l-4$ after taking a
difference. A zero-slope binomial is only constant and cannot affect
the coefficient at degree $l-3$. Summing, with $\sum_i c_i=1$,
proves~\eqref{eq:centralleading}.

For $l=3$, all singleton subset sums are below $W/2$ and all pair
sums are above it. The full product contributes $-2(s-2)$, each
two-factor product has difference $-1$, and their coefficients sum
to one. This gives $-2s+3$ exactly.

Finally, the degree bound for $Q_w$ follows either from
\eqref{eq:centralleading} or the exact cancellation
\eqref{eq:topsymmetry}. Lemma~\ref{lem:righttail} gives $Q_w(t)\le0$
at all integers of the indicated parity where it represents the
central difference. A nonzero polynomial with this property has
negative leading coefficient.
\end{proof}

For a nonzero rational polynomial $P(t)=\sum_{j=0}^d p_jt^j$, define
\[
 b(P)=1+\left\lceil\frac{\sum_{j<d}|p_j|}{|p_d|}\right\rceil,
 \qquad b(0)=1.
\]
For every real $t\ge b(P)$, a nonzero $P(t)$ has the sign of $p_d$.
Indeed, if $d\ge1$ and $t\ge1$, the absolute sum of its lower terms
is at most $t^{d-1}\sum_{j<d}|p_j|<|p_d|t^d$.
The constant case is immediate.

\begin{proof}[Proof of Theorem~\ref{thm:dilation}, including an effective threshold]
If $2\max_iw_i\ge W$, split a dilated word at its maximum $M$ and
write $D$ for the remaining weight. Here $M\ge D$. First-half
increase reaches rank $D$, whereas the maximum identity gives
\[
 \Del F_{XMY}(n)=-[z^{n-M-1}]F_XF_Y\le0
 \qquad(n\ge D+2).
\]
Only the difference at $D+1$ is undecided, so take $T(w)=1$.
Otherwise $l\ge3$ and the central curvature is negative.

Fix a parity of $t$, and write $N=tW$, $L=\lceil N/2\rceil$,
$V=\lfloor(N+l-1)/2\rfloor$, $U=\lceil(N+l-1)/2\rceil$.
The offsets $n-N/2$ in the following list depend only on the parity:
\[
 L<n\le V;\qquad n=V;\qquad n=U+1.
\]
For the first two groups, $s<l/2$, so $P_s$ has positive leading
coefficient. For the last index, $s>l/2$, so the leading coefficient
is negative. Include $Q_w$ as well when $N+l$ is even.
Take $T(w)$ to be the maximum of $4l+4$ and all the integers $b(P)$
for these finitely many polynomials, over both parities.
This is a finite exact-arithmetic construction.

For $t\ge T(w)$, Lemma~\ref{lem:half} gives nonnegative differences
through $L$; the selected positive polynomials give positive
differences from $L+1$ through $V$, as well as at $V$ itself.
Lemma~\ref{lem:righttail} gives nonpositive differences at every
index after $V$. Hence the sequence is unimodal and $V$ is a mode.
The strict difference at $V$ excludes earlier modes, and the strict
negative difference at $U+1$ excludes later modes.
If $U=V$, this proves uniqueness. Otherwise the difference at $U$
is $Q_w(t)$; it is either strictly negative beyond its bound or
identically zero. This proves the precise eventual alternatives.
\end{proof}

For example, the supplied exact procedure gives $T(w)=32$ for
$w=(1,4,2,5,3,2,4)$, with unique eventual mode
$\lfloor(21t+6)/2\rfloor$ in both parities. This is a sufficient
threshold, not a claim of the least possible threshold.
The procedure computes and certifies central polynomials; it does not
test finitely many dilations and extrapolate.

\begin{corollary}[Mean rank]\label{cor:dilationmean}
Choose a distinct smaller composition uniformly from the ideal of
$tw$, including the empty composition, and let its weight be $K_t$.
For every fixed nonempty $w$,
\[
 \mathbb E K_t=\frac{tW+l-1}{2}+O(t^{-1}).
\]
\end{corollary}
\begin{proof}
A product $\prod_{i\in I}S_{tw_i}$ has mass
$t^{|I|}\prod_{i\in I}w_i$ and mean exponent
$(t\sum_{i\in I}w_i+|I|)/2$. The full polynomial has mass
$t^l\prod_iw_i+O(t^{l-1})$.
Subtract $(tW+l)/2$ times this mass from its first moment.
The terms missing exactly one factor contribute
\[
 -\tfrac12t^l\prod_iw_i\sum_i c_i+O(t^{l-1})
 =-\tfrac12t^l\prod_iw_i+O(t^{l-1});
\]
all other proper terms contribute $O(t^{l-1})$.
Divide by the total mass. The same calculation holds for $l=1$,
when the sole proper term is the empty product.
\end{proof}
